\section*{Solution}

\subsection*{Part (b) — Width of $M_3$/$M_4$ and Mixer Conversion Gain}

% -------------------------------------------------------
\paragraph{What is an IF component?}

This circuit is a heterodyne receiver.  $M_1$ amplifies the RF signal at
$f_\text{RF} = 5.20\,\text{GHz}$.  The mixer ($M_2$, $M_3$, $M_4$) then
\emph{down-converts} it to the intermediate frequency
\[
  f_\text{IF} = f_\text{RF} - f_\text{LO} = 5.20 - 5.13 = 70\,\text{MHz}.
\]
When a product $i_{d2}(t) = G_{m2}\,V_\text{rf}\cos(\omega_\text{RF}t)$ is
multiplied by the LO switching function $s(t)$ at $\omega_\text{LO}$, the
output spectrum contains components at
$|m\,f_\text{LO} \pm n\,f_\text{RF}|$ for non-negative integers $m,n$.
The \textbf{IF component} is the spectral line at $f_\text{IF} = 70\,\text{MHz}$
--- the down-converted replica of the RF signal that all subsequent baseband
processing acts on.  All other mixing products (at $f_\text{RF}+f_\text{LO}$,
$3f_\text{LO}\pm f_\text{RF}$, etc.) are filtered by the IF selectivity.

% -------------------------------------------------------
\paragraph{How the switching mixer produces the IF component.}

$M_2$ acts as a \emph{transconductor}: it converts the RF voltage at its gate
to a tail drain current $i_\text{tail}(t) = G_{m2}\,v_\text{rf}(t)$.
$M_3$ and $M_4$ form a \emph{current-commutating pair}: driven
differentially by $V_\text{LO}$, they steer $i_\text{tail}$ alternately
between the two 900\,$\Omega$ load branches at the LO rate.  When the LO
amplitude is large enough to switch $M_{3}$/$M_4$ \emph{completely} every
half-cycle (\textbf{hard switching}), the differential steering function is
a unit-amplitude square wave:
\[
  s(t) = \frac{4}{\pi}\!\left[\cos(\omega_\text{LO}t)
         - \tfrac{1}{3}\cos(3\omega_\text{LO}t) + \cdots\right].
\]
The differential IF output current is $\Delta i_\text{out} = s(t)\cdot i_\text{tail}(t)$.
Picking only the term at $\omega_\text{IF} = \omega_\text{RF}-\omega_\text{LO}$:
\[
  \Delta I_\text{IF} = \frac{2}{\pi}\,G_{m2}\,V_\text{rf},
  \qquad\Rightarrow\qquad
  V_\text{IF} = \frac{2}{\pi}\,G_{m2}\,V_\text{rf}\cdot R_L,
\]
where $R_L = 900\,\Omega$.  With $V_\text{rf} = V_\text{in}/2$
(the RF voltage at $M_2$'s gate referenced to the source), the
\textbf{voltage conversion gain of the mixer} is
\begin{equation}\label{eq:CG}
  \boxed{C_V = \frac{V_\text{IF}}{V_\text{in}/2}
             = \frac{2}{\pi}\,G_{m2}\,R_L.}
\end{equation}
The factor $2/\pi\approx 0.64$ is the fundamental Fourier coefficient of the
square-wave switching function and is the \emph{maximum possible} switching
efficiency.  It is only realized under hard switching.

% -------------------------------------------------------
\paragraph{Why the width of $M_3$/$M_4$ sets the conversion gain.}

$M_3$ and $M_4$ are a differential pair with tail current $I_{SS} = 0.15\,\text{mA}$
supplied by $M_2$.  At their quiescent point each carries $I_{SS}/2$, so
their overdrive voltage is
\[
  V_\text{GS,eff,Q}
    = \sqrt{\frac{I_{SS}}{\mu_n C_\text{ox}\,(W_{34}/L)}}.
\]
For a square-law differential pair the complete-switching condition is
(derivable from $\Delta I_D = 0$ for the off transistor):
\begin{equation}\label{eq:hs}
  V_\text{LO,peak} \;\geq\; \sqrt{2}\,V_\text{GS,eff,Q}
    = \sqrt{\frac{2\,I_{SS}}{\mu_n C_\text{ox}\,(W_{34}/L)}}.
\end{equation}
\begin{itemize}
  \item \textbf{If $W_{34}$ is too small}: $V_\text{GS,eff,Q}$ is large,
    condition~\eqref{eq:hs} fails, $M_3$/$M_4$ never fully switch off, and
    some tail current recirculates rather than being steered to the output.
    The effective switching factor drops below $2/\pi$ and $C_V$ is
    \emph{below its maximum}.
  \item \textbf{If $W_{34}$ satisfies condition~\eqref{eq:hs}}: hard
    switching is achieved and $C_V = (2/\pi)\,G_{m2}\,R_L$, its maximum.
    Making $W_{34}$ larger than this threshold does not further increase
    $C_V$ (it only adds unnecessary parasitic capacitance at the IF output
    node).
\end{itemize}
Therefore, \textbf{maximum conversion gain is achieved at the minimum $W_{34}$
that satisfies the hard-switching condition} with the given LO amplitude.

% -------------------------------------------------------
\paragraph{Finding $W_{34}$.}

Set $V_\text{LO,peak} = \sqrt{2}\,V_\text{GS,eff,Q}$ and solve for $W_{34}$:
\[
  V_\text{LO,peak}^2 = \frac{2\,I_{SS}\,L}{\mu_n C_\text{ox}\,W_{34}}
  \quad\Longrightarrow\quad
  \boxed{W_{34} = \frac{2\,I_{SS}\,L}{\mu_n C_\text{ox}\,V_\text{LO,peak}^2}.}
\]
The LO is specified as a peak-to-peak \emph{differential} amplitude of
$0.9\,\text{V}$, so the differential peak is
$V_\text{LO,peak} = 0.9/2 = 0.45\,\text{V}$.
Substituting $I_{SS} = 0.15\,\text{mA}$, $L = 0.18\,\mu\text{m}$, and
$\mu_n C_\text{ox} = 270\,\mu\text{A/V}^2$:
\[
  W_{34}
    = \frac{2\times0.15\times10^{-3}\times0.18\times10^{-6}}
           {270\times10^{-6}\times(0.45)^2}
    = \frac{5.4\times10^{-11}}{5.47\times10^{-8}}
    \approx \mathbf{1\,\mu\text{m}}.
\]

% -------------------------------------------------------
\paragraph{Conversion gain under this condition.}

With hard switching established, equation~\eqref{eq:CG} applies directly.
$G_{m2}$ is the effective RF transconductance of $M_2$.  Its source is
degenerated by the parallel combination of $10\,\text{k}\Omega$ and
$0.25\,\text{pF}$; at $f_\text{RF} = 5.2\,\text{GHz}$ the capacitor dominates
($Z_{0.25\,\text{pF}} \approx 122\,\Omega \ll 10\,\text{k}\Omega$), so
\[
  G_{m2} = \frac{g_{m2}}{1 + g_{m2}\,Z_s},
  \quad Z_s = \frac{1}{j\,2\pi\,f_\text{RF}\times 0.25\,\text{pF}}
              \approx -j\,122\,\Omega,
\]
with the intrinsic transconductance
\[
  g_{m2} = \sqrt{2\,\mu_n C_\text{ox}\,\frac{W_2}{L}\,I_{D2}}
          = \sqrt{2\times270\times10^{-6}\times\frac{W_2}{0.18\times10^{-6}}
                  \times0.15\times10^{-3}},
\]
where $W_2$ is obtained from part~(a).  The voltage conversion gain is then
\[
  \boxed{C_V = \frac{2}{\pi}\,
               \frac{g_{m2}}{\left|1+g_{m2}\cdot(-j122\,\Omega)\right|}
               \times 900\,\Omega.}
\]
If source degeneration at RF is neglected ($g_{m2}\times122\,\Omega\ll1$),
this simplifies to $C_V \approx (2/\pi)\,g_{m2}\times900$.
